\part{Understanding}

% ─────────────────────────────────────────────────────────────────
\chapter{Celestial coordinates}
\label{ch:coords}

\section{The celestial sphere}

Seen from the ground, the night sky behaves like a sphere turning from east to
west about an axis through the celestial poles. Stars lie at wildly different
distances, but their position on that fictitious sphere is all one needs to
find and follow them.

The Earth completes one turn in \textbf{23~h 56~min 4~s} --- the
\emph{sidereal day}, shorter than the 24 hours of the solar day. That, and not
the clock, is the rate a mount must turn at.

\section{Right ascension}

\textbf{Right ascension} (RA) is the celestial equivalent of longitude. It is
counted in \textbf{hours} rather than degrees, from 0~h to 24~h, starting at
the vernal point.

\begin{formula}
\[
1\,\text{h} = 15\degre \qquad 1\,\text{min} = 15' \qquad 1\,\text{s} = 15''
\]
So RA = 02~h 44~min 08~s is $2\times15\degre + 44\times0.25\degre +
8\times\frac{15}{3600}\degre = 41.03\degre$.
\end{formula}

\begin{info}
RA is the critical coordinate for tracking: a one-second timing error is worth
15 arcseconds on the sky at the celestial equator. At 1.6\arcsec/pixel that is
about ten pixels of trailing.
\end{info}

\section{Declination}

\textbf{Declination} (DEC) is the equivalent of latitude: the angular distance
north or south of the celestial equator, from $-90\degre$ to $+90\degre$. On a
correctly aligned mount the declination axis does not move while tracking.

\section{Hour angle}

The \textbf{hour angle} measures how far an object stands from the local
meridian. It grows continuously as the Earth turns.

\begin{formula}
\[
\text{HA} = \text{local sidereal time} - \text{RA}
\]
An RA tolerance can be expressed in arcseconds or in seconds of time. The
conversion depends on declination:
\[
\text{seconds of time} = \frac{\text{arcseconds}}{15 \cos(\text{DEC})}
\]
\end{formula}

\begin{tip}
Near the celestial pole an RA error barely moves the star on the sensor: the
cosine of the declination damps it. Circumpolar objects are therefore more
forgiving of RA tracking errors --- but not of declination ones.
\end{tip}

% ─────────────────────────────────────────────────────────────────
\chapter{Equatorial mounts}
\label{ch:mounts}

\section{Principle}

An \textbf{equatorial mount} cancels the Earth's rotation by turning the tube
about an axis parallel to the Earth's own: the polar axis.

\begin{itemize}
  \item \textbf{Polar axis (RA)}: turns at the sidereal rate.
  \item \textbf{Declination axis}: still while tracking, in theory.
\end{itemize}

\begin{formula}
\[
\omega = \frac{1\,296\,000''}{86\,164.09\,\text{s}} \approx
15.041\,\text{arcseconds per second}
\]
\end{formula}

Any departure from that rate, and any misalignment of the polar axis, becomes a
drift the image records.

\section{10Micron mounts}

10Micron mounts (GM1000 HPS and above) carry \textbf{absolute encoders}: they
know where they point even after a power cut, where a relative encoder loses
it. To that they add a worm drive, periodic-error correction and, above all, an
internal \textbf{pointing model} built from several dozen stars.

That model is what makes \textbf{unguided} imaging possible: the mount does not
merely turn at the sidereal rate, it applies corrected rates that absorb
flexure, refraction and the residual polar misalignment.

\begin{warning}
A direct consequence, often misunderstood: on a modelled mount the
\textbf{polar alignment error lives inside the model} and is already
compensated by the tracking rates. A residual declination drift therefore does
not call for the azimuth screw, but for \textbf{more model points near the
target}.
\end{warning}

These mounts speak an \textbf{extended LX200}: beyond the standard commands
they expose raw axis positions, applied rates, model state and environmental
data. MountMonitor uses those extensions where available.

% ─────────────────────────────────────────────────────────────────
\chapter{What degrades tracking}
\label{ch:errors}

No mount tracks perfectly. The causes overlap, and each leaves a different
signature in the record.

\section{Periodic error}

A worm is never perfectly round nor perfectly centred. Its defect repeats with
every turn and produces an oscillation of fixed period --- a few minutes, most
often. That is exactly what a Fourier analysis can pull out
(chapter~\ref{ch:fft}).

\section{Polar alignment}

If the polar axis does not point exactly at the celestial pole, the star drifts
slowly in declination, the faster the further one is from the meridian. On a
mount without a model the answer is mechanical. On a modelled mount, see the
warning in the previous chapter.

\section{Atmospheric refraction}

The atmosphere lifts objects the more they are low: some $34'$ at the horizon,
$1'$ at $45\degre$ of altitude. Refraction therefore changes through the night,
and its drift is easily mistaken for a misalignment.

\section{Mechanical disturbance and vibration}

Wind, a cable pulling, ground movement, a meridian flip: so many one-off jolts,
none of them periodic. They show as isolated spikes, and that is what tells
them apart.

\section{Commanded moves}

These are not errors at all. Between exposures the sequencer deliberately moves
the mount --- \textbf{dither} to decorrelate noise, \textbf{re-centering} when
the target has slipped --- and the mount \emph{stays} at its new position.
Confusing them with tracking error is the costliest mistake one can make
reading a record, and chapter~\ref{ch:stats} is devoted to it.

% ─────────────────────────────────────────────────────────────────
\chapter{Reading the numbers}
\label{ch:stats}

\begin{warning}
This is the most important chapter in the manual. A tracking record is read
wrongly with disconcerting ease, and the figure that comes out is then wrong by
a factor of several hundred.
\end{warning}

\section{The signal is a staircase}

Here is what a night of unguided imaging actually records: the mount holds its
position to a fraction of an arcsecond, then the sequencer dithers it, and it
holds its new position just as well. The trace is not a noisy line about a
mean: it is a \textbf{staircase}.

A standard deviation taken over the whole night then measures the height of the
steps, not the quality of the tracking. And removing a straight line --- what a
classic detrending does --- does not remove a staircase.

\begin{info}
On a real eight-and-a-half hour session the raw combined RMS came out at
\textbf{86.3\arcsec}, a damning verdict. The mount was in fact holding each
position to \textbf{0.18\arcsec}. The two figures differ by a factor of 480.
\end{info}

\section{Jitter}

\textbf{Jitter} is the spread the mount shows \emph{while} it holds a position.
It is the only figure that smears a star on the sensor, and the one
MountMonitor bases its rating on.

It is measured between repositionings, never about a straight line. The program
cuts the record at the steps, drops two samples after each --- the move itself
is not tracking --- and takes the median of the per-step spreads.

\begin{tip}
This approach has a property worth noting: a sequence that dithers on every
frame does not come out worse than one that never does. The figure describes
the mount, not the sequencer.
\end{tip}

\section{Commanded moves}

They are counted and classified separately:

\begin{center}
\rowcolors{2}{rowalt}{white}
\begin{tabularx}{\textwidth}{lX}
\toprule
\textbf{Class} & \textbf{What it is} \\
\midrule
Dither & A small move between exposures, to decorrelate noise. \\
Re-centering & A larger move, bringing the target back where it belongs. \\
Unsettled & The spread that follows stays several times the usual value. That
is the one worth looking at. \\
\bottomrule
\end{tabularx}
\end{center}

\begin{warning}
A threshold crossing is not a move. A 70\arcsec{} re-centering takes several
samples to complete and trips the threshold at each one: one night read as 98
crossings for 53 actual moves, with a median amplitude ten times too small.
MountMonitor groups crossings less than ten seconds apart.
\end{warning}

\section{Slow drift}

It is quoted per hour and per exposure. On an unguided mount every dither
cancels it: what matters is not the drift over the night, but the displacement
it produces \emph{during one exposure}.

\begin{formula}
\[
\text{displacement over one exposure} = \text{drift (\arcsec/h)} \times
\frac{\text{exposure (s)}}{3600}
\]
A drift of $-16.8\arcsec$/h therefore moves the star by $0.84\arcsec$ over a
180~s exposure.
\end{formula}

\section{Running standard deviation}

Computed live over a 60 to 900 second window, it serves to see \emph{when}
something happened, not to grade the night. After a target change, an
acquisition gap or a repositioning, its window straddles the move and describes
that rather than the tracking: those values are excluded from the report.

\section{The rating}

It rests on the combined jitter, with thresholds suited to the type of mount.
An unguided 10Micron is judged on wider thresholds than a guided mount, because
the orders of magnitude are not the same.

\begin{center}
\rowcolors{2}{rowalt}{white}
\begin{tabularx}{\textwidth}{lXX}
\toprule
\textbf{Rating} & \textbf{Guided mount} & \textbf{Unguided precision} \\
\midrule
Excellent & $< 0.5\arcsec$ & $< 2.0\arcsec$ \\
Good      & $< 1.5\arcsec$ & $< 4.0\arcsec$ \\
Fair      & $< 3.0\arcsec$ & $< 8.0\arcsec$ \\
\bottomrule
\end{tabularx}
\end{center}

\begin{tip}
Always relate the rating to your sampling. At 1.6\arcsec/pixel, a
0.5\arcsec{} jitter shakes the star by a third of a pixel: invisible on the
final image.
\end{tip}

% ─────────────────────────────────────────────────────────────────
\chapter{Fourier analysis}
\label{ch:fft}

\section{What it is for}

A Fourier transform breaks a signal into a sum of sinusoids. On a tracking
record it brings out what \emph{repeats} --- and a worm periodic error repeats
by construction.

MountMonitor looks in the \textbf{2 to 15 minute} band, where the worm periods
of common mounts fall.

\section{Two precautions}

\begin{itemize}
  \item The transform runs on \textbf{detrended} deviations. A linear drift is
  not periodic, but over a finite window it produces a large low-frequency
  component that would hide what is being looked for.
  \item The \textbf{Nyquist limit} requires sampling at least twice per period.
  At 2~Hz, periods below one second are out of reach --- no real inconvenience,
  no mechanical error lives there.
\end{itemize}

\section{Reading a spectrum}

A sharp peak at a stable period marks a periodic error. A flat spectrum means
there is no detectable one --- the case of 10Micron mounts after correction.
Broad, wandering peaks are noise, usually wind.
